\subsection*{Sprint 2: Despliegue automático}
\addcontentsline{toc}{subsection}{Sprint 2: Despliegue automático}

El segundo sprint se centra en automatizar el proceso de despliegue del sistema
\textbf{Muncher} en la nube, estableciendo un entorno de desarrollo accesible y
reproducible. El objetivo principal es lograr que, con cada cambio en la rama
\texttt{main}, el sistema se despliegue automáticamente en un entorno remoto,
garantizando una integración continua entre el desarrollo local y la
infraestructura en la nube.

Para asegurar la consistencia y trazabilidad del proyecto, la infraestructura
se define íntegramente como código (\textit{Infrastructure as Code}, IaC).

El despliegue automático debe ejecutarse de manera inmediata como respuesta a
los \textit{commits} realizados en la rama principal, eliminando pasos manuales
y reduciendo el riesgo de errores humanos durante la actualización del entorno
de desarrollo.

\subsubsection*{Objetivos iniciales}

\begin{itemize}
	\item Configurar un entorno de desarrollo en la nube accesible para el equipo.
	\item Definir la infraestructura como código utilizando herramientas de IaC.
	\item Implementar un flujo de integración y despliegue continuo (CI/CD) que se active
	      automáticamente ante cambios en la rama \texttt{main}.
	\item Garantizar la reproducibilidad del entorno para futuras fases del proyecto,
	      como el despliegue en producción.
\end{itemize}

\subsubsection*{Historias de usuario completadas}

\begin{center}
	\begin{tabularx}{\textwidth}{|l|X|c|}
		\hline
		\textbf{Identificador} & \textbf{Nombre}                                                 & \textbf{Puntos de historia} \\
		\hline
		\usref{US-DE-01}       & Despliegue en la nube del front-end en el entorno de desarollo  & 8                           \\
		\hline
		\usref{US-DE-02}       & Despliegue en la nube del front-end en el entorno de producción & 2                           \\
		\hline
	\end{tabularx}
\end{center}

\subsubsection*{Retrospectiva}

El sprint concluyó con la implementación exitosa de la infraestructura de nube
del proyecto \textbf{Muncher}. El desafío principal se centró en el análisis
comparativo de los distintos proveedores de nube y en la selección de una
plataforma que cumpliera los requisitos técnicos del sistema sin comprometer su
flexibilidad ni aumentar los costes de operación. Este proceso exigió una
evaluación profunda de configurabilidad, compatibilidad con contenedores,
servicios gestionados y políticas de precios. Las plataformas ligeras
resultaron atractivas por su sencillez operativa y por la amplitud de sus capas
gratuitas, pero el análisis detallado de dichas capas reveló que la persistencia
de datos no podía sostenerse en ellas más allá de un plazo breve, lo que
habría obligado a rehacer el despliegue sobre otro proveedor poco después de
tener el sistema en funcionamiento. Por este motivo se seleccionó
\textbf{Amazon Web Services}, que cubre los requisitos del sistema sin
condicionar su continuidad y permite construir el entorno de desarrollo y el de
producción sobre la misma infraestructura.

Una vez seleccionada la plataforma, el despliegue se abordó definiendo la
infraestructura como código. El sistema se describió mediante una plantilla
común para ambos entornos —desarrollo y producción—, garantizando consistencia
en la infraestructura y simplificando la replicación de recursos: la misma
definición se instancia tantas veces como entornos sean necesarios, con las
diferencias justas entre ellos expresadas como parámetros. La aplicación web
estática se desplegó automáticamente en el entorno de desarrollo con cada
commit a la rama \texttt{main}, integrándose con el flujo de integración
continua del proyecto. Los despliegues a producción, en cambio, no se
desencadenan a partir de ningún cambio en el código: se lanzan de forma
explícita, aunque el proceso que ejecutan está igualmente automatizado y es el
mismo que se emplea en desarrollo. Lo que se reserva a una decisión humana es
únicamente el momento de desplegar, de modo que se preserva la estabilidad del
servicio sin renunciar a la reproducibilidad del procedimiento.

La restricción de acceso al entorno de desarrollo, recogida en
\reqref{RNF-SE-05}, se resolvió durante este mismo sprint. Un entorno de
desarrollo desplegado en la nube es, por defecto, tan accesible como el de
producción: cualquiera que conozca su dirección puede visitarlo, cuando su
finalidad es servir únicamente al equipo que lo utiliza para verificar los
cambios. Para evitarlo se interpuso una autenticación básica en el punto de
entrada de la aplicación.

El sprint permitió consolidar las bases técnicas de la infraestructura en la
nube, culminando con un entorno operativo estable, reproducible y de bajo
coste. Además, la experiencia reafirmó la eficacia de definir la infraestructura
como código y demostró que una configuración bien estructurada puede ofrecer una
transición transparente entre entornos de desarrollo y producción sin
complejidad adicional.

Las URL de los entornos desplegados son:

\begin{itemize}
	\item Entorno de desarrollo: \url{https://dev.muncher-cook.com}
	\item Entorno de producción: \url{https://muncher-cook.com}
\end{itemize}
